\documentclass[10pt,a4paper]{amsart}
\usepackage[margin=19mm]{geometry}
\usepackage{amsmath,amssymb,mathtools}
\usepackage[scr=rsfs,cal=euler]{mathalfa}
\usepackage{xfrac}
\usepackage[unicode,hidelinks,bookmarks=false]{hyperref}
\newcommand{\ie}{i.e.\ }
\newcommand{\cf}{cf.\ }
\newcommand{\resp}{respectively\ }
\newcommand{\Subsection}{\subsection}
\providecommand{\nobreakdash}{\nobreak\hbox{-}}
\setlength{\emergencystretch}{2em}
\multlinegap0pt
\usepackage{amssymb}
\newtheorem{proposition}{Proposition}
\title{\textsf{Closing in on the kernel of an operator between Banach spaces}}
\author{Douglas S. Bridges}
\date{Source: arXiv:2604.26259v2 (CC BY 4.0)}
\begin{document}
\maketitle

\noindent\emph{Source and licence.} \textsf{Closing in on the kernel of an operator between Banach spaces}. Version arXiv:2604.26259v2 (CC BY 4.0), distributed by arXiv under the Creative Commons Attribution 4.0 International licence, http://creativecommons.org/licenses/by/4.0/, as recorded in the arXiv licence metadata for this exact version and shown on the abstract page https://arxiv.org/abs/2604.26259v2. Full licence notice: \texttt{licenses/arxiv-2604.26259-CC-BY-4.0.txt}.\par\smallskip

\noindent\textit{Editorial scope.} Counted unit: the article's proposition apr20p1 (source0.tex lines 203--236). The article's complete original proof of it is delivered with it (source0.tex lines 238--302, one \texttt{proof} environment of 65 lines). No mathematics is written, repaired or re-derived: the statement and the proof are the article's own lines, in the article's own notation, and no retained line is substituted or re-flowed.  No further source-local context is needed: every cross-reference of every retained line resolves inside the two delivered spans.  Nothing was omitted from any retained span, so the package carries no omission record.  No retained line contains a citation, so the wrapper carries no bibliography and no bibliography entry.  No retained line contains a figure environment, an image-inclusion command or drawing code, so no image is delivered and the package declares no asset dependency.  Editorial formatting only, in addition: an AMS \texttt{amsart} preamble that restates the article's own declarations and macros used by the retained lines, namely the environments \texttt{proposition} on separate counters, together with the standard packages those lines need.  The retained environments print in this document as Proposition 1 in order of appearance, whereas the article prints the same items with the numbering of its own full document.  The complete native source of the article as distributed in the arXiv e-print for this exact version is bundled unchanged as \texttt{sources/199375/source0.tex}.
\begin{proposition}
\label{apr20p1}Let $T$ be a nonzero linear mapping of a normed space $X$ into
a normed space $Y$. Then the following conditions are equivalent.

\begin{itemize}
\item[\emph{(a)}] There exists $r>0$ such that%
\[
\delta_{r}\equiv\inf\left\{  \left\Vert x\right\Vert :x\in X,\left\Vert
Tx\right\Vert \geq r\right\}
\]
exists and is positive;

\item[\emph{(b)}] $T$ is normable,\footnote{%
%TCIMACRO{\TeXButton{sf}{\normalfont\sf}}%
%BeginExpansion
\normalfont\sf
%EndExpansion
\textsf{Recall that a linear mapping }$T:X\rightarrow Y$\textsf{\ between
normed spaces is \textbf{normable} if its \textbf{operator norm} }$\left\Vert
T\right\Vert \equiv\sup\left\{  \left\Vert Tx\right\Vert :x\in X,\left\Vert
x\right\Vert \leq1\right\}  $\textsf{\ exists.}} with positive norm;

\item[\emph{(c)}] $\delta_{r}$ exists for each $r>0$.
\end{itemize}

%

%TCIMACRO{\TeXButton{noindent}{\noindent}}%
%BeginExpansion
\noindent
%EndExpansion
If any, and therefore all, of these conditions holds, then $\left\Vert
T\right\Vert =r\delta_{r}^{-1}$ for each $r>0$.
\end{proposition}

\begin{proof}
Suppose that $r$ and $\delta_{r}\ $exist as in (i). If $x\in X$ and
$\left\Vert x\right\Vert <1$, then $\left\Vert T(\delta_{r}x)\right\Vert
\not \geq r$, so $\left\Vert Tx\right\Vert =\left\Vert T(\delta_{r}%
x)\right\Vert \delta_{r}^{-1}\leq r\delta_{r}^{-1}$. It readily follows that
$T$ is a bounded linear mapping with bound $r\delta_{r}^{-1}$. To prove that
$r\delta_{r}^{-1}$ is the norm of $T$, let $0<\varepsilon<r\delta_{r}^{-1}$.
Then $0<1-r^{-1}\delta_{r}\varepsilon<1$ and therefore%
\[
\frac{\delta_{r}}{1-r^{-1}\delta_{r}\varepsilon}>\delta_{r}.
\]
Choose, in turn, $t$ such that%
\[
0<t<\frac{\delta_{r}}{1-r^{-1}\delta_{r}\varepsilon}-\delta_{r}%
\]
and $\xi\in X$ such that $\left\Vert T\xi\right\Vert \geq r$ and $\left\Vert
\xi\right\Vert <\delta_{r}+t$. With%
\[
x=\frac{r}{\left(  \delta_{r}+t\right)  \left\Vert T\xi\right\Vert }\xi
\]
we have%
\[
\left\Vert x\right\Vert =\frac{r}{\left\Vert T\xi\right\Vert }\cdot
\frac{\left\Vert \xi\right\Vert }{\delta_{r}+t}<1
\]
and%
\[
\left\Vert Tx\right\Vert =\frac{r}{\delta_{r}+t}>\frac{r(1-r^{-1}\delta
_{r}\varepsilon)}{\delta_{r}}=r\delta_{r}^{-1}-\varepsilon.
\]
Since $\varepsilon>0$ is arbitrary, it follows that%
\[
\left\Vert T^{-1}\right\Vert =\sup\left\{  \left\Vert x\right\Vert :x\in
X,\left\Vert Tx\right\Vert \leq1\right\}
\]
exists and equals $r\delta_{r}^{-1}$.

Now suppose that $T$ is normable and $\left\Vert T\right\Vert >0$. For each
$r>0$, if $\left\Vert Tx\right\Vert \geq r$, then $\left\Vert x\right\Vert
\geq r\left\Vert T\right\Vert ^{-1}$. On the other hand, given $\varepsilon
>0$, pick, in turn, $s$ such that $0<s<\left\Vert T\right\Vert $ and%
\[
\frac{1}{\left\Vert T\right\Vert -s}<\frac{1}{\left\Vert T\right\Vert }%
+r^{-1}\varepsilon
\]
and $\xi\in X$ with $\left\Vert \xi\right\Vert =1$ and $\left\Vert
T\xi\right\Vert >\left\Vert T\right\Vert -s$. Setting%
\[
x=\frac{r}{\left\Vert T\right\Vert -s}\xi,
\]
we have%
\[
\left\Vert Tx\right\Vert =\frac{r\left\Vert T\xi\right\Vert }{\left\Vert
T\right\Vert -s}>r
\]
and%
\[
\left\Vert x\right\Vert =\frac{r}{\left\Vert T\right\Vert -s}<r\left(
\left\Vert T\right\Vert ^{-1}+r^{-1}\varepsilon\right)  =r\left\Vert
T\right\Vert ^{-1}+\varepsilon\text{.}%
\]
Since $r$ and $\varepsilon$ are arbitrary, we conclude that for each $r>0$,
$\delta_{r}$ exists and equals $r\left\Vert T\right\Vert ^{-1}$. Thus (b)
implies (c), which trivially implies (a).
\end{proof}

\end{document}
